\section{\faWrench\ Multimodalità\\\small{Text-to-speech}} % (fold)
\label{sec:spring-tts}
%
\begin{frame}[t,fragile] \frametitle{Multimodalità}
    \framesubtitle{Text-to-speech}
    {\footnotesize
    \begin{itemize}[leftmargin=10pt,align=right]
        \onslide<1->\item[\alert{\faArrowCircleRight}] Sintetizzazione sonora a partire da testo (opposto della trascrizione)
        \onslide<2->\item[\alert{\faExclamationTriangle}] Non possibile utilizzare le capacità multimodali dei LLM generalisti
        \onslide<3->\item[\alert{\faArrowCircleRight}] Compatibilità Spring AI (al 11/2025)
        \begin{itemize}[leftmargin=10pt,align=right]
            \item[\alert{\faArrowCircleRight}] OpenAI API
            \item[\alert{\faArrowCircleRight}] ElevenLabs
        \end{itemize}
        \onslide<4->\item[\alert{\faArrowCircleRight}] Possibilità di scelta fra tono, voci, velocità, \ldots
        \onslide<5->\item[\alert{\faArrowCircleRight}] Ambiti applicativi
        \begin{itemize}[leftmargin=10pt,align=right]
            \item[\alert{\faArrowCircleRight}] IVR (Interactive Voice Response) - Sistema telefonico automatizzato di seconda generazione
            \item[\alert{\faArrowCircleRight}] Tecnologie assistive decicate alle diverse abilità
            \item[\alert{\faArrowCircleRight}] \ldots
        \end{itemize}        
    \end{itemize}
    }
\end{frame}   
%
\begin{frame}[t,fragile] \frametitle{Progetto Spring AI}
    \framesubtitle{Text-to-speech: Spring AI}
    	\vspace*{-.7cm}
        \begin{block}{Interfaccia di base}
			{\tiny\inputminted{java}{code/TextToSpeechModel.java}}
    	\end{block}
        \begin{itemize}[leftmargin=10pt,align=right]
            \onslide<1->\item[\alert{\faExclamationTriangle}] \textit{Breaking change} al passaggio verso Spring AI 1.1.0
            \onslide<2->\item[\alert{\faExclamationTriangle}] Transizione verso entità funzionali e \textit{major refactoring}
            \onslide<3->\item[\alert{\faExclamationTriangle}] Consultare la documentazione costantemente
        \end{itemize}         
\end{frame}
%
\begin{frame}[t,fragile] \frametitle{Progetto Spring AI}
    \framesubtitle{Text-to-speech: OpenAI con impostazioni di \textit{default}}
    	\vspace*{-.7cm}
        \begin{block}{\textit{Controller} di sintetizzazione audio}
			{\tiny\inputminted{java}{code/AudioController.java}}
    	\end{block}
        \begin{itemize}[leftmargin=10pt,align=right]
            \onslide<1->\item[\alert{\faArrowCircleRight}] Impostazioni di \textit{default}
            \onslide<2->\begin{itemize}[leftmargin=10pt,align=right]
                \item[\alert{\faArrowCircleRight}] GPT 4.0 mini TTS
                \item[\alert{\faArrowCircleRight}] \textit{Response format} in JSON
                \item[\alert{\faArrowCircleRight}] Velocità di sintetizzazione 1.0
                \item[\alert{\faArrowCircleRight}] Voce Alloy
            \end{itemize}
            \onslide<3->\item[\alert{\faArrowCircleRight}] Istanziato un \textit{bean} \alert{\texttt{OpenAiAudioSpeechModel}}
        \end{itemize}      
\end{frame}
%
\begin{frame}[t,fragile] \frametitle{Progetto Spring AI}
    \framesubtitle{Text-to-speech: OpenAI con impostazioni \textit{custom} dinamiche}
    	\vspace*{-.7cm}
        \begin{block}{\textit{Controller} di sintetizzazione audio}
			{\tiny\inputminted{java}{code/AudioControllerCustom.java}}
    	\end{block}
        \begin{itemize}[leftmargin=10pt,align=right]
            \onslide<1->\item[\alert{\faArrowCircleRight}] Popolato un \alert{\texttt{TextToSpeechPrompt}} con tutte le sovrascritture volute rispetto al \textit{default}     
            \onslide<2->\item[\alert{\faArrowCircleRight}] \alert{\texttt{speechResponse.getResult().getOutput()}} ripulisce l'output da tutti i metadati restituendo il dato grezzo (\texttt{byte[]})
        \end{itemize}         
\end{frame}
%
\begin{frame}[t,fragile] \frametitle{Multimodalità}
\framesubtitle{Text-to-speech: OpenAI \textit{compatible} API}
    {\footnotesize
    \begin{itemize}[leftmargin=10pt,align=right]
        \onslide<1->\item[\alert{\faExclamationTriangle}] Impossibile utilizzare compatibilità per LLM generalista
        \onslide<2->\item[\alert{\faExclamationTriangle}] Obbligatorio \textit{workaround} di \textit{call} diretta API del AI \textit{provider} tramite Spring Web
    \end{itemize}
    \begin{block}{Approccio REST per sintetizzazione vocale: Gemini - singolo}
		{\tiny\inputminted{bash}{code/gemini-tts-single.sh}}
    \end{block}
    }
\end{frame} 
%
\begin{frame}[t,fragile] \frametitle{Multimodalità}
\framesubtitle{Text-to-speech: OpenAI \textit{compatible} API}
    \vspace*{-.7cm}
    {\footnotesize
    \begin{block}{Approccio REST per sintetizzazione vocale: Gemini - conversazione}
		{\tiny\inputminted{bash}{code/gemini-tts-conversation.sh}}
    \end{block}
    }
\end{frame}    
%
\begin{frame}[t,fragile] \frametitle{Progetto Spring AI}
    \framesubtitle{Applicazione e passaggi}
    {\small
    \begin{itemize}[leftmargin=10pt,align=right]
        \onslide<1->\item[\alert{\faArrowCircleRight}] Servizio \textit{text-to-speech} Gemini
        \begin{itemize}[leftmargin=10pt,align=right]
            \onslide<2->\item[\alertedcircled{1}] Modifica dipendenze di progetto
            \onslide<3->\item[\alertedcircled{2}] Creazione modelli Gemini TTS \textit{request} e \textit{response}
            \onslide<4->\item[\alertedcircled{3}] Creazione modello per informazioni voci disponibili
            \onslide<5->\item[\alertedcircled{4}] Creazione convertitore \textit{output} audio Gemini a formato riproducibile nel \textit{web}
            \onslide<6->\item[\alertedcircled{5}] Modifica controllore MVC
            \onslide<7->\item[\alertedcircled{6}] \textit{Test} delle funzionalità con Postman/Insomnia
        \end{itemize}
    \end{itemize}
    }
\end{frame}
%
\begin{frame}[t,fragile] \frametitle{Progetto Spring AI}
    \framesubtitle{Text-to-speech: Gemini}
    {\footnotesize
        \begin{block}{Dipendenze applicativo}
			{\tiny\inputminted{xml}{code/pom.xml}}
    	\end{block} 
    }
\end{frame}
%
\begin{frame}[t,fragile] \frametitle{Progetto Spring AI}
    \framesubtitle{Text-to-speech: Gemini}
        \begin{block}{Modello richiesta Gemini TTS}
			{\tiny\inputminted{java}{code/GeminiTtsRequest.java}}
    	\end{block}
\end{frame}
%
\begin{frame}[t,fragile] \frametitle{Progetto Spring AI}
    \framesubtitle{Text-to-speech: Gemini}
        \begin{block}{Modello risposta Gemini TTS}
			{\tiny\inputminted{java}{code/GeminiTtsResponse.java}}
    	\end{block}
\end{frame}
%
\begin{frame}[t,fragile] \frametitle{Progetto Spring AI}
    \framesubtitle{Text-to-speech: Gemini}
    	\vspace*{-.7cm}
        \begin{block}{Utilità convertitore PCM -> WAV}
			{\tiny\inputminted{java}{code/PcmToWavConverter.java}}
    	\end{block}
\end{frame}
%
\begin{frame}[t,fragile] \frametitle{Progetto Spring AI}
    \framesubtitle{Text-to-speech: Gemini}
        \begin{block}{Modello informazioni voci disponibili}
			{\tiny\inputminted{java}{code/VoiceInfo.java}}
    	\end{block}
\end{frame}
%
\begin{frame}[t,fragile] \frametitle{Progetto Spring AI}
    \framesubtitle{Text-to-speech: Gemini}
    	\vspace*{-.7cm}
        \begin{block}{Implementazione controllore REST}
			{\tiny\inputminted{java}{code/GeminiMultimodalityController.java}}
    	\end{block}
\end{frame}
%
\begin{frame}[t,fragile] \frametitle{Progetto Spring AI}
    \framesubtitle{Text-to-speech: Gemini}
    	\vspace*{-.7cm}
        \begin{block}{Pagina web}
			{\tiny\inputminted{html}{code/gemini-tts.html}}
    	\end{block}
\end{frame}
%
\begin{frame}[fragile] \frametitle{Codice}
    \framesubtitle{Branch di riferimento}
	\begin{center}
		{\scriptsize \href{https://github.com/simonescannapieco/spring-ai-advanced-dgroove-venis-code.git}{\texttt{https://github.com/simonescannapieco/spring-ai-advanced-dgroove-venis-code.git}}}\\
		\textit{Branch:} \alert{\texttt{18-spring-ai-gemini-ollama-multimodality-tts}}
	\end{center}
\end{frame}